\Opensolutionfile{ans}[ans/ansTN-PhanDangLuu1204]

\noindent \textbf{PHẦN I. Câu trắc nghiệm nhiều phương án lựa chọn (3,0 điểm). Học sinh trả lời từ câu 1 đến câu 12. Mỗi câu hỏi thí sinh chỉ chọn một phương án.}
\vspace{0.2cm}

\begin{ex}	\macau{ID: VL12-PDL1204-C01}Hệ thức chuyển đổi từ nhiệt độ dóng theo thang nhiệt độ Celsius sang thang nhiệt độ tuyệt đối Kelvin là
	\choice
	{\True $T(\text{K})=t(^\circ\text{C})+273$.}
	{$T(\text{K})=t(^\circ\text{C})-273$.}
	{$T(\text{K})=1,8t(^\circ\text{C})-273$.}
	{$T(\text{K})=1,8t(^\circ\text{C})+273$.}
	\loigiai{
		Mối liên hệ giá trị nhiệt độ giữa thang Celsius và thang Kelvin tuân theo phương trình dịch mốc chuẩn mực là: $T(\text{K}) = t(^\circ\text{C}) + 273,15$. Lấy làm tròn số nguyên ta được $T = t + 273$.
	}
\end{ex}

\begin{ex}	\macau{ID: VL12-PDL1204-C02}Khi nói về các đặc trưng cấu trúc của chất ở thể rắn, phát biểu nào sau đây là đúng?
	\choice
	{Chất rắn có thể tích riêng nhưng không có hình dạng xác định, dễ nén hơn vật ở thể lỏng.}
	{Chất rắn luôn chiếm toàn bộ thể tích bình chứa và có thể nén được dễ dàng.}
	{\True Chất rắn có thể tích và hình dạng riêng xác định, rất khó bị nén.}
	{Chất rắn có hình dạng và thể tích riêng xác định, rất dễ bị nén.}
	\loigiai{
		Ở thể rắn, các lực tương tác liên phân tử là vô cùng mạnh mẽ giữ chặt các phân tử dao động nhiệt quanh vị trí cân bằng cố định, giúp chất rắn giữ được hình dạng và thể tích riêng xác định, đồng thời rất khó bị nén ép cơ học.
	}
\end{ex}

\begin{ex}	\macau{ID: VL12-PDL1204-C03}Hệ thức nào dưới đây là phù hợp với quá trình một khối khí lí tưởng chứa trong bình kín có thể tích không đổi bị nung nóng từ bên ngoài? (Bỏ qua sự giãn nở vì nhiệt của bình).
	\choice
	{\True $\Delta U=Q$ khi $Q>0$.}
	{$\Delta U=Q$ khi $Q<0$.}
	{$\Delta U=A$ khi $A>0$.}
	{$\Delta U=A$ khi $A<0$.}
	\loigiai{
		Vì khối khí chứa trong bình kín có thể tích không đổi ($\Delta V = 0$) nên khối khí hoàn toàn không thực hiện công và không nhận công từ bên ngoài ($A = 0$). Khi bị nung nóng, khối khí nhận nhiệt lượng nên $Q > 0$. Áp dụng Nguyên lí I ta được: $\Delta U = Q$ với $Q > 0$.
	}
\end{ex}

\begin{ex}	\macau{ID: VL12-PDL1204-C04}Công thức nào sau đây mô tả đúng định luật của Nguyên lí I nhiệt động lực học?
	\choice
	{$\Delta U=A-Q$.}
	{$\Delta U=Q-A$.}
	{$A=\Delta U+Q$.}
	{\True $\Delta U=A+Q$.}
	\loigiai{
		Biểu thức chuẩn tắc của Nguyên lí I nhiệt động lực học khẳng định độ biến thiên nội năng của hệ bằng tổng công và nhiệt lượng hệ trao đổi với môi trường ngoài: $\Delta U = A + Q$.
	}
\end{ex}

\begin{ex}	\macau{ID: VL12-PDL1204-C05}Những ngày mùa hè thời tiết nóng bức, khi ta dùng nước được lấy lên từ giếng khoan sâu thì thấy nước rất mát. Ngược lại, những ngày mùa đông khi nhiệt độ không khí xuống thấp thì dùng nước giếng lại thấy ấm. Nguyên nhân chính của hiện tượng này là vì
	\choice
	{nước cách nhiệt tốt nên hoàn toàn không bị ảnh hưởng bởi thời tiết bên trên.}
	{nước giếng khoan là mạch nước ngầm chảy từ vùng núi cao nên chịu ảnh hưởng thời tiết.}
	{nước giếng cách mặt đất một khoảng xa nên không bị ảnh hưởng bởi thời tiết.}
	{\True nước có nhiệt dung riêng lớn hơn đất và không khí nên độ thay đổi nhiệt độ của nước mạch ngầm sâu là rất nhỏ so với không khí.}
	\loigiai{
		Nước lỏng có hằng số nhiệt dung riêng lớn hơn rất nhiều so với đất đá và không khí. Do đó, mạch nước ngầm sâu dưới lòng đất giữ nhiệt độ tương đối ổn định quanh năm. Mùa hè không khí ngoài trời nóng hơn nước giếng nên ta cảm thấy mát; mùa đông không khí lạnh hơn nước giếng nên ta cảm thấy ấm.
	}
\end{ex}

\begin{ex}	\macau{ID: VL12-PDL1204-C06}Khi trời lạnh, ô tô bật điều hòa sưởi ấm và đóng kín tất cả các cửa kính, hành khách ngồi trên ô tô sau một thời gian sẽ quan sát thấy hiện tượng vật lí gì xuất hiện?
	\choice
	{Không có hiện tượng gì đặc biệt xuất hiện.}
	{\True Hơi nước ngưng tụ tạo thành các giọt nước bám ở mặt phía trong của kính xe.}
	{Nước bốc hơi dữ dội trên bề mặt xe.}
	{Hơi nước ngưng tụ tạo thành các giọt nước bám ở mặt phía ngoài của kính xe.}
	\loigiai{
		Hành khách ngồi trong xe kín thở ra hơi ẩm chứa nhiều hơi nước ấm. Khi hơi nước ấm này di chuyển va chạm vào bề mặt phía trong của lớp kính xe (đang tiếp xúc trực tiếp với không khí lạnh bên ngoài nên có nhiệt độ thấp), hơi nước bị mất nhiệt và ngưng tụ tạo thành màng sương giọt nước ở phía bên trong kính.
	}
\end{ex}

\begin{ex}	\macau{ID: VL12-PDL1204-C07}Số đo đại lượng hằng số nhiệt hóa hơi riêng $L$ của chất lỏng là thông tin kĩ thuật vô cùng cần thiết trực tiếp phục vụ cho việc thiết kế thiết bị nào sau đây?
	\choice
	{\True quạt điều hòa hơi nước.}
	{máy biến áp nguồn.}
	{máy điều hòa không khí chạy lốc.}
	{nhiệt kế lưỡng kim.}
	\loigiai{
		Quạt hơi nước vận hành làm mát dựa trên cơ chế thu nhiệt từ không khí để hóa hơi hoàn toàn các giọt nước lỏng mịn phun ra, do đó thông số nhiệt hóa hơi riêng của nước là cốt lõi để tính toán công suất lưu lượng nước cần thiết.
	}
\end{ex}

\begin{ex}	\macau{ID: VL12-PDL1204-C08}Định nghĩa nào sau đây đúng khi nói về đại lượng hằng số nhiệt nóng chảy riêng của một chất?
	\choice
	{là toàn bộ nhiệt lượng cần cung cấp cho vật để làm vật nóng chảy.}
	{Anhiệt lượng cần để làm cho một đơn vị khối lượng chất đó nóng chảy một phần lẻ.}
	{là trị số nhiệt độ nóng chảy riêng biệt của khối chất rắn kết tinh.}
	{\True là nhiệt lượng cần thiết để làm cho một đơn vị khối lượng ($1\text{ kg}$) chất đó nóng chảy hoàn toàn ở nhiệt độ nóng chảy mà không làm thay đổi nhiệt độ.}
	\loigiai{
		Nhiệt nóng chảy riêng $\lambda$ của một chất là nhiệt lượng cần cung cấp cho đúng $1\text{ kg}$ chất đó chuyển pha nóng chảy hoàn toàn ở nhiệt độ nóng chảy xác định dưới áp suất cho trước.
	}
\end{ex}

\begin{ex}	\macau{ID: VL12-PDL1204-C09}Cách xác định vạch các mốc nhiệt độ tiêu chuẩn trong hệ nhiệt giai bách phân Celsius là
	\choice
	{Lấy nhiệt độ của nước khi đóng băng là ($100^\circ\text{C}$) và nhiệt độ sôi của nước làm chuẩn ($10^\circ\text{C}$).}
	{Lấy nhiệt độ của nước khi đóng băng là ($10^\circ\text{C}$) và nhiệt độ sôi của nước làm chuẩn ($0^\circ\text{C}$).}
	{\True Lấy nhiệt độ của nước khi đóng băng là ($0^\circ\text{C}$) và nhiệt độ sôi của nước làm chuẩn ($100^\circ\text{C}$).}
	{Lấy nhiệt độ của nước khi đóng băng là ($10^\circ\text{C}$) and nhiệt độ sôi của nước làm chuẩn ($100^\circ\text{C}$).}
	\loigiai{
		Thang nhiệt độ Celsius chọn điểm đóng băng của nước tinh khiết làm vạch mốc $0^\circ\text{C}$ và điểm hơi nước đang sôi bão hòa dưới áp suất tiêu chuẩn làm vạch mốc $100^\circ\text{C}$.
	}
\end{ex}

\begin{ex}	\macau{ID: VL12-PDL1204-C10}Tìm điểm tương đồng, giống nhau cốt lõi giữa hai thang đo nhiệt độ Celsius và thang đo nhiệt độ tuyệt đối Kelvin.
	\choice
	{Mọi giá trị nhiệt độ đo được trên cả hai thang đều mang giá trị đại số dương.}
	{Có chung một vạch mốc nhiệt độ trùng với nhiệt độ điểm ba của nước tinh khiết.}
	{\True Độ rộng của một vạch chia độ ($\Delta t = 1^\circ\text{C}$) trên thang này bằng đúng độ rộng của một vạch chia ($\Delta T = 1\text{ K}$) trên thang kia.}
	{Có chung điểm mốc nhiệt độ xuất phát đặt tại điểm “Độ không tuyệt đối".}
	\loigiai{
		Khoảng cách nhiệt độ giữa điểm băng tan và điểm nước sôi ở cả hai thang đều được chia nhỏ thành đúng 100 phần bằng nhau, do đó khoảng chênh lệch nhiệt độ thay đổi $1^\circ\text{C}$ hoàn toàn trùng khít bằng với lượng thay đổi $1\text{ K}$.
	}
\end{ex}

\begin{ex}	\macau{ID: VL12-PDL1204-C11}Với các kí hiệu chuẩn đã quy ước trong sách giáo khoa về nhiệt lượng, biểu thức công thức nào dưới đây \textbf{không} phải là công thức dùng để tính nhiệt lượng mà vật thu vào hay tỏa ra khi thay đổi nhiệt độ?
	\choice
	{\True $Q=m c\left(t_{2}+t_{1}\right)$.}
	{$Q=m c\left(T_{2}-T_{1}\right)$.}
	{$Q=m c \Delta T$.}
	{$Q=m c \Delta t$.}
	\loigiai{
		Nhiệt lượng tỏa ra hay thu vào khi thay đổi nhiệt độ tỉ lệ thuận với biến thiên nhiệt độ (hiệu số nhiệt độ), do đó công thức viết tổng hai nhiệt độ $t_1 + t_2$ ở phương án A là sai.
	}
\end{ex}

\begin{ex}	\macau{ID: VL12-PDL1204-C12}Trong các hiện tượng tự nhiên và đời sống sau, hiện tượng nào có liên quan trực tiếp đến tiến trình chuyển pha của sự nóng chảy?
	\choice
	{Đun nóng một nồi nước lỏng trên bếp điện.}
	{Đốt cháy bấc ngọn đèn dầu hỏa.}
	{Cho một cốc nước lỏng vào ngăn đá tủ lạnh.}
	{\True Thả cục nước đá rắn vào cốc nước lỏng ấm.}
	\loigiai{
		Khi thả cục nước đá rắn vào cốc nước ấm, cục nước đá thu nhiệt lượng từ nước và chuyển pha trạng thái từ thể rắn sang thể lỏng, đây chính là hiện tượng nóng chảy.
	}
\end{ex}

\Closesolutionfile{ans}

%%%%%%%%%%%%%%%%%%%%%%%%%%%%%%%%%%%%%%%%%%%%%%%%%%%%%%%%%%%%%%%%%%%%%%%%%
% PHẦN II: CÂU TRẮC NGHIỆM ĐÚNG SAI
%%%%%%%%%%%%%%%%%%%%%%%%%%%%%%%%%%%%%%%%%%%%%%%%%%%%%%%%%%%%%%%%%%%%%%%%%
\noindent \textbf{PHẦN II. Câu trắc nghiệm đúng sai (2,0 điểm). Học sinh trả lời từ câu 1 đến câu 2. Trong mỗi ý a), b), c), d) ở mỗi câu, thí sinh chọn đúng hoặc sai.}
\Opensolutionfile{ans}[ans/ansTF-PhanDangLuu1204]

\begin{ex}	\macau{ID: VL12-PDL1204-C13}Người ta dùng một chiếc ấm điện hoạt động với công suất định mức $P = 1000\text{ W}$ để đun sôi một lượng nước có khối lượng $m=1,2\text{ kg}$ từ nhiệt độ ban đầu $t_{0}=20^\circ\text{C}$ cho đến khi sôi hoàn toàn ở $t=100^\circ\text{C}$ dưới áp suất khí quyển tiêu chuẩn. Cho biết hằng số nhiệt dung riêng và ẩn nhiệt hóa hơi riêng của nước lần lượt là $c=4,2 \cdot 10^{3}\text{ J/(kg.K)}$ và $L=2,26 \cdot 10^{6}\text{ J/kg}$.
	\choiceTF[t]
	{Giả sử ấm điện không thất thoát nhiệt, thời gian tối thiểu để đun lượng nước trên sôi phải là 6 phút 20 giây.}
	{\True Biết thực tế hiệu suất của ấm nước đạt $90\%$. Sau khi nước sôi, người này quên rút nguồn và để ấm nước tiếp tục sôi thêm đúng $1\text{ phút}$, khối lượng nước lỏng bị hóa hơi bốc thoát đi xấp xỉ bằng $23,9\text{ g}$.}
	{Nhiệt lượng mà lượng nước thu vào để thực hiện tiến trình sôi lượng chất ở nhiệt độ cố định $100^\circ\text{C}$ được tính bằng biểu thức công thức $Q=L \cdot m$.}
	{\True Ý nghĩa của thông số hiệu suất ấm nước đạt $90\%$ có nghĩa là chỉ có đúng $90\%$ lượng điện năng năng lượng tiêu thụ được biến đổi thành nhiệt năng hữu ích truyền vào đun nước.}
	\loigiai{
		\begin{itemchoice}
			\itemch Nhiệt lượng có ích cần thiết để nâng nhiệt khối nước đến điểm sôi: 
			$Q_{\text{ích}} = m \cdot c \cdot (t - t_0) = 1,2 \cdot 4200 \cdot (100 - 20) = 403200\text{ J}$. \\
			Thời gian lý tưởng tối thiểu: $\tau_{\text{min}} = \frac{Q_{\text{ích}}}{P} = \frac{403200}{1000} = 403,2\text{ s} = 6\text{ phút } 43,2\text{ giây}$ (không phải 6 phút 20 giây).
			\itemch Trong thời gian $\tau = 1\text{ phút} = 60\text{ s}$ tiếp theo, năng lượng điện bếp tỏa ra là $W = P \cdot \tau = 1000 \cdot 60 = 60000\text{ J}$. \\
			Nhiệt lượng thực tế truyền vào phục vụ hóa hơi nước dưới hiệu suất $90\%$: $Q_{\text{hh}} = H \cdot W = 0,9 \cdot 60000 = 54000\text{ J}$. \\
			Khối lượng nước lỏng bị hóa hơi bốc đi: $\Delta m = \frac{Q_{\text{hh}}}{L} = \frac{54000}{2,26 \cdot 10^6} \approx 0,02389\text{ kg} \approx 23,9\text{ g}$.
			\itemch Công thức $Q = L \cdot m$ dùng để tính nhiệt lượng cung cấp cho khối lượng $m$ chất lỏng hóa hơi hoàn toàn ở nhiệt độ sôi, không phải nhiệt lượng để nâng nước đến khi bắt đầu sôi.
			\itemch Theo định nghĩa hiệu suất thiết bị nhiệt, hiệu suất phản ánh tỉ lệ phần trăm năng lượng có ích nhận được trên tổng năng lượng tiêu thụ toàn phần.
		\end{itemchoice}
	}
\end{ex}

\begin{ex}	\macau{ID: VL12-PDL1204-C14}Khi người ta thực hiện truyền một lượng nhiệt lượng có độ lớn bằng $Q = 9 \cdot 10^{3}\text{ J}$ cho khối khí lí tưởng chứa bên trong một bình xi-lanh hình trụ kín, khối khí dãn nở đẩy tấm pít-tông dịch chuyển tịnh tiến lên trên một đoạn dài $10\text{ cm}$. Biết tiết diện ngang của lòng xi-lanh dóng bằng $0,05\text{ m}^2$, lượng nội năng đại số của chất khí được tăng thêm một lượng bằng $\Delta U = 5 \cdot 10^{3}\text{ J}$ và coi áp suất của khí không đổi trong suốt quá trình biến đổi.
	\choiceTF[t]
	{\True Trị số áp suất nội tại không đổi của khối chất khí đẩy pít-tông bằng $8 \cdot 10^{5}\text{ N/m}^2$.}
	{\True Tổng hợp quá trình nhiệt động trên, chất khí đồng thời nhận nhiệt lượng và thực hiện sinh công làm gia tăng biến thiên nội năng của hệ.}
	{Độ biến thiên nội năng tính theo quy ước dấu Nguyên lí I biểu diễn bởi hệ thức $\Delta U = A + Q$ với các giá trị thành phần thỏa mãn dấu: $A < 0$, $Q > 0$ và $\Delta U < 0$.}
	{\True Lượng thể tích không gian hình trụ mà khối khí đã giãn nở mở rộng thêm bằng đúng $0,005\text{ m}^3$.}
	\loigiai{
		\begin{itemchoice}
			\itemch Khí nhận nhiệt $Q = +9 \cdot 10^3\text{ J}$ và tăng nội năng $\Delta U = +5 \cdot 10^3\text{ J}$. \\
			Theo Nguyên lí I, công đại số khí nhận vào bằng: $A = \Delta U - Q = 5 \cdot 10^3 - 9 \cdot 10^3 = -4 \cdot 10^3\text{ J}$. \\
			Công khí thực hiện sinh ra đẩy pít-tông: $A' = |A| = 4000\text{ J}$. \\
			Độ dời hành trình pít-tông: $s = 10\text{ cm} = 0,1\text{ m}$. Lực đẩy pít-tông của áp suất khí: $F = \frac{A'}{s} = \frac{4000}{0,1} = 40000\text{ N}$. \\
			Trị số áp suất khí: $p = \frac{F}{S} = \frac{40000}{0,05} = 800000\text{ N/m}^2 = 8 \cdot 10^5\text{ Pa}$.
			\itemch Khí nhận nhiệt lượng từ ngoài ($Q>0$) đồng thời dãn nở đẩy pít-tông sinh công lực hướng ra ngoài ngoại cảnh ($A<0$) làm tăng nội năng.
			\itemch Biểu thức Nguyên lí I viết dạng $\Delta U = A + Q$ là đúng, các dấu thành phần $A < 0$ và $Q > 0$ đúng, nhưng trị số biến thiên nội năng đề cho tăng nên dấu đại số phải mang trị số dương: $\Delta U > 0$.
			\itemch Lượng thể tích hình trụ tăng thêm do pít-tông dịch chuyển: $\Delta V = S \cdot s = 0,05\text{ m}^2 \cdot 0,1\text{ m} = 0,005\text{ m}^3$.
		\end{itemchoice}
	}
\end{ex}

\Closesolutionfile{ans}

%%%%%%%%%%%%%%%%%%%%%%%%%%%%%%%%%%%%%%%%%%%%%%%%%%%%%%%%%%%%%%%%%%%%%%%%%
% PHẦN III: CÂU TRẮC NGHIỆM TRẢ LỜI NGShort
%%%%%%%%%%%%%%%%%%%%%%%%%%%%%%%%%%%%%%%%%%%%%%%%%%%%%%%%%%%%%%%%%%%%%%%%%
\noindent \textbf{PHẦN III. Câu trắc nghiệm trả lời ngắn (2,0 điểm). Học sinh trả lời từ câu 1 đến câu 4.}
\Opensolutionfile{ans}[ans/ansSA-PhanDangLuu1204]

\begin{ex}	\macau{ID: VL12-PDL1204-C15}Để tiêu diệt các bào tử nấm độc hại bám ngoài vỏ vỏ và kích thích nhanh tiến trình nảy mầm tự nhiên của hạt giống lúa, người nông dân thường sử dụng một kinh nghiệm dân gian lâu đời là pha nước tắm hạt theo tỉ lệ “hai sôi, ba lạnh”. Tức là nước ấm pha trộn được tạo ra bằng cách hòa chung 2 phần thể tích nước sôi ở vạch $100^\circ\text{C}$ với 3 phần thể tích nước lạnh. Hãy xác định trị số nhiệt độ của khối nước ấm pha thu được theo đơn vị độ Celsius ($^\circ\text{C}$), biết nhiệt độ của nguồn nước lạnh ban đầu đem pha bằng $20^\circ\text{C}$ và coi hệ lý tưởng.

	\par\shortans[oly]{52}
	\loigiai{
		Do khối lượng lỏng tỉ lệ thuận trực tiếp với thể tích, ta coi khối lượng nước sôi hòa vào là $m_1 = 2$ phần và khối lượng nước lạnh là $m_2 = 3$ phần. Lập phương trình cân bằng nhiệt lượng thu tỏa của hệ trộn:
		$$Q_{\text{tỏa}} = Q_{\text{thu}} \Rightarrow m_1 \cdot c \cdot (t_{\text{sôi}} - t_{\text{cb}}) = m_2 \cdot c \cdot (t_{\text{cb}} - t_{\text{lạnh}})$$
		$$2 \cdot (100 - t_{\text{cb}}) = 3 \cdot (t_{\text{cb}} - 20) \Rightarrow 200 - 2t_{\text{cb}} = 3t_{\text{cb}} - 60$$
		$$260 = 5t_{\text{cb}} \Rightarrow t_{\text{cb}} = \frac{260}{5} = 52^\circ\text{C.}$$
	}
\end{ex}

\begin{ex}	\macau{ID: VL12-PDL1204-C16}Người ta tiến hành thực hiện một công cơ học có độ lớn bằng $100\text{ J}$ tác dụng lên một hệ vật chất, nhưng đồng thời trong suốt chu trình biến đổi lượng nhiệt lượng bị thất thoát truyền cản ra môi trường không khí ngoài đo được có độ lớn bằng $60\text{ J}$. Hỏi trị số lượng nội năng của hệ vật chất trên đã biến thiên một lượng bằng bao nhiêu Jun?

	\par\shortans[oly]{40}
	\loigiai{
		Trích xuất các giá trị đại số tuân theo quy ước dấu Nguyên lí I của nhiệt động lực học:
		- Hệ nhận công cơ học từ ngoại lực tác dụng $\Rightarrow A = +100\text{ J}$.
		- Hệ tỏa nhiệt lượng (thất thoát nhiệt) ra ngoài môi trường $\Rightarrow Q = -60\text{ J}$.
		Độ biến thiên nội năng tổng hợp của hệ vật:
		$$\Delta U = Q + A = -60 + 100 = +40\text{ J.}$$
	}
\end{ex}

\begin{ex}	\macau{ID: VL12-PDL1204-C17}Lịch sử khí tượng thế giới từng ghi nhận một sự biến đổi nhiệt độ không khí siêu đột ngột diễn ra tại thị trấn Spearfish, bang South Dakota (Hoa Kỳ) vào ngày 22/01/1943. Lúc 7h30 sáng, nhiệt độ ngoài trời đo được là $-20^\circ\text{C}$, nhưng chỉ sau đúng 2 phút ngắn ngủi kế tiếp, nhiệt độ đã dâng tăng vọt lên mức $7,2^\circ\text{C}$. Xác định độ tăng nhiệt độ trung bình trong khoảng thời gian 2 phút đột biến đó dóng theo đơn vị đo Kelvin trên phút ($\text{K/phút}$)? (Kết quả làm tròn lấy đến một chữ số thập phân sau dấu phẩy).

	\par\shortans[oly]{13,6}
	\loigiai{
		- Độ biến thiên tăng nhiệt độ đo theo thang Celsius: $\Delta t = t_2 - t_1 = 7,2 - (-20) = 27,2^\circ\text{C}$. \\
		- Vì độ rộng của một độ chia vạch thang Celsius bằng đúng độ rộng một phân chia thang Kelvin, lượng tăng nhiệt tuyệt đối đạt: $\Delta T = \Delta t = 27,2\text{ K}$. \\
		- Tốc độ gia tăng nhiệt độ trung bình tính trong khoảng thời gian $\Delta \tau = 2\text{ phút}$ là:
		$$\text{Tốc độ tăng} = \frac{\Delta T}{\Delta \tau} = \frac{27,2\text{ K}}{2\text{ phút}} = 13,6\text{ K/phút.}$$
	}
\end{ex}

\begin{ex}	\macau{ID: VL12-PDL1204-C18}Một thỏi nhôm đặc có khối lượng cơ học $m = 1,0\text{ kg}$ đang ở nhiệt độ phòng ổn định $25^\circ\text{C}$. Biết các thông số kĩ thuật: ẩn nhiệt nóng chảy riêng của nhôm bằng $\lambda = 4 \cdot 10^5\text{ J/kg}$, hằng số nhiệt dung riêng bằng $880\text{ J/(kg.K)}$ và điểm nóng chảy của nhôm ở $658^\circ\text{C}$. Tính tổng nhiệt lượng cần thiết phải cung cấp để đưa thỏi nhôm trên hóa lỏng hoàn toàn ở đúng nhiệt độ nóng chảy của nó theo đơn vị đo kilôjun ($\text{kJ}$)? (Kết quả chỉ lấy phần số nguyên).

	\par\shortans[oly]{957}
	\loigiai{
		Tiến trình thu nhiệt lượng toàn phần trải qua hai giai đoạn liên tiếp:
		- Chặng 1: Thu nhiệt nâng nhiệt độ khối nhôm từ điểm phòng $25^\circ\text{C}$ lên đến thềm điểm nóng chảy $658^\circ\text{C}$:
		$$Q_1 = m \cdot c \cdot (t_{\text{nc}} - t_0) = 1,0 \cdot 880 \cdot (658 - 25) = 880 \cdot 633 = 557040\text{ J.}$$
		- Chặng 2: Thu nhiệt phá vỡ mạng tinh thể để hóa lỏng hoàn toàn tại vạch điểm $658^\circ\text{C}$:
		$$Q_2 = m \cdot \lambda = 1,0 \cdot 4 \cdot 10^5 = 400000\text{ J.}$$
		Tổng nhiệt lượng toàn phần cung cấp cho hệ nhôm:
		$$Q = Q_1 + Q_2 = 557040 + 400000 = 957040\text{ J} = 957,04\text{ kJ.}$$
		Lấy phần số nguyên theo đúng yêu cầu ta thu được giá trị chỉ số bằng $957\text{ kJ}$.
	}
\end{ex}

\Closesolutionfile{ans}

%%%%%%%%%%%%%%%%%%%%%%%%%%%%%%%%%%%%%%%%%%%%%%%%%%%%%%%%%%%%%%%%%%%%%%%%%
% PHẦN IV: ĐỀ BÀI TỰ LUẬN TÍNH TOÁN (3 ĐIỂM)
%%%%%%%%%%%%%%%%%%%%%%%%%%%%%%%%%%%%%%%%%%%%%%%%%%%%%%%%%%%%%%%%%%%%%%%%%
\noindent \textbf{PHẦN IV. Tự luận (3,0 điểm). Học sinh trả lời từ câu 1 đến câu 3.}
\vspace{0.2cm}

\begin{ex}	\macau{ID: VL12-PDL1204-C19}Người ta thực hiện truyền cấp một nhiệt lượng có độ lớn bằng $25\text{ J}$ cho một khối khí lí tưởng chứa bên trong một bình xi-lanh đặt nằm ngang. Khối chất khí hấp thụ năng lượng dãn nở sinh lực đẩy cơ học thực hiện đẩy tấm pít-tông trượt dời chỗ trong lòng xi-lanh được một khoảng hành trình dài $10\text{ cm}$. Biết lực ma sát trượt cản trở xuất hiện ở bề mặt tiếp xúc giữa pít-tông và vách thành xi-lanh có độ lớn bằng $20\text{ N}$ và coi pít-tông trượt chuyển động đều. Hãy tính trị số độ biến thiên nội năng $\Delta U$ đại số của khối khí khí?
	\loigiai{
		- Quy đổi khoảng hành trình dịch chuyển dời chỗ sang mét chuẩn hệ SI: $s = 10\text{ cm} = 0,1\text{ m}$. \\
		- Vì tấm pít-tông thực hiện trượt tịnh tiến thẳng đều, độ lớn lực đẩy do áp suất chất khí sinh ra cân bằng hoàn toàn đối với độ lớn lực ma sát trượt cản phá bề mặt: $F_{\text{đẩy}} = F_{\text{ms}} = 20\text{ N}$. \\
		- Độ lớn công cơ học do khối chất khí thực hiện sinh ra giải phóng hướng ra ngoại cảnh bằng:
		$$A' = F_{\text{đẩy}} \cdot s = 20 \cdot 0,1 = 2\text{ J.}$$
		- Hệ chất khí thực hiện sinh công nên hằng số công đại số nhận dấu âm quy ước dấu: $A = -2\text{ J}$. \\
		- Trích xuất nhiệt lượng chất khí nhận vào từ nguồn ngoài: $Q = +25\text{ J}$. \\
		- Áp dụng hệ thức định luật toán học Nguyên lí I nhiệt động lực học để tính độ biến thiên nội năng:
		$$\Delta U = Q + A = 25 + (-2) = +23\text{ J.}$$
		Giá trị dương cho biết nội năng chất khí tăng một lượng bằng $23\text{ J}$.
	}
\end{ex}

\begin{ex}	\macau{ID: VL12-PDL1204-C20}Một chiếc khay làm bằng sắt có khối lượng cơ học $1,2\text{ kg}$ được bọc lớp cách điện và làm nóng bằng một lò sưởi điện có công suất định mức bằng $500\text{ W}$ hoạt động chạy liên tục trong vòng $4\text{ phút}$. Người ta dùng thiết bị cảm biến ghi nhận thấy nhiệt độ của khay sắt đã gia tăng đều đặn dâng từ vạch mức đầu $22^\circ\text{C}$ đạt lên đến vạch thềm cuối $45^\circ\text{C}$. Giả định hệ kín lý tưởng và bỏ qua hoàn toàn mọi sự thất thoát hao phí nhiệt lượng truyền ra ngoài môi trường không khí. Hãy xác định giá trị hằng số nhiệt dung riêng của vật liệu sắt từ phép đo thực nghiệm trên?
	\loigiai{
		- Quy đổi thông số khoảng thời gian đun của máy sưởi sang giây chuẩn hệ SI:
		$$\tau = 4\text{ phút} = 4 \cdot 60 = 240\text{ s.}$$
		- Tổng năng lượng điện năng do máy sưởi điện tiêu thụ tỏa ra cung cấp cho hệ bằng:
		$$W = P \cdot \tau = 500\text{ W} \cdot 240\text{ s} = 120000\text{ J.}$$
		- Do hệ kín lý tưởng không có hao phí thoát nhiệt, toàn bộ điện năng biến đổi trọn vẹn thành nhiệt lượng hữu ích truyền cấp nâng nhiệt khay sắt: $Q = W = 120000\text{ J}$. \\
		- Áp dụng biểu thức công thức tính lượng nhiệt nâng nhiệt độ đưa vào hệ vật chất:
		$$Q = m \cdot c \cdot (t_2 - t_1) \Rightarrow 120000 = 1,2 \cdot c \cdot (45 - 22)$$
		$$120000 = 1,2 \cdot c \cdot 23 \Rightarrow 120000 = 27,6 \cdot c$$
		$$\Rightarrow c = \frac{120000}{27,6} \approx 4347,8\text{ J/(kg.K).}$$
	}
\end{ex}

\begin{ex}	\macau{ID: VL12-PDL1204-C21}Người ta bỏ một cục nước đá lạnh có khối lượng $m$ đang ở mốc nhiệt độ đóng băng $0^\circ\text{C}$ vào trong một chiếc xô nhựa đang chứa nước lỏng. Tổng khối lượng hỗn hợp cả nước và đá chứa trong xô đạt giá trị $M = 10\text{ kg}$. Tiến hành dùng thiết bị theo dõi đo đạc thu được đồ thị mô tả sự biến thiên của nhiệt độ $t(^\circ\text{C})$ phụ thuộc dóng theo trục thời gian $\tau\text{ (phút)}$ như hình vẽ đính kèm bên cạnh. Cho biết hằng số nhiệt dung riêng của nước lỏng bằng $c = 4200\text{ J/(kg.K)}$, ẩn nhiệt nóng chảy riêng của nước đá bằng $\lambda = 3,4 \cdot 10^5\text{ J/kg}$, công suất cung cấp nhiệt lượng đun của môi trường là hằng hằng số không đổi và hệ kín cô lập hoàn toàn. Hãy tính toán xác định hằng số lượng khối lượng $m$ của cục nước đá ban đầu bằng bao nhiêu kilôgam (kg)?
	\loigiai{
		- Khảo sát đồ thị thực nghiệm chuyển pha dóng theo thời gian để phân tích bản chất:
		\begin{itemize}
			\item \textbf{Giai đoạn 1 (từ $\tau = 0 \rightarrow 50\text{ phút}$):} Đường đồ thị nằm ngang bão hòa ở thềm mức cố định $0^\circ\text{C}$. Đây là phân đoạn hệ nhận nhiệt lượng từ nguồn ngoài chỉ để phục vụ riêng cho tiến trình chuyển pha nóng chảy tan rã hoàn toàn khối lượng $m$ của cục nước đá rắn. Khoảng thời gian nóng chảy kéo dài: $\Delta \tau_1 = 50 - 0 = 50\text{ phút}$. Lượng nhiệt cung cấp chặng này: $Q_1 = P \cdot \Delta \tau_1 = \lambda \cdot m$.
			\item \textbf{Giai đoạn 2 (từ $\tau = 50 \rightarrow 60\text{ phút}$):} Đường đồ thị là một đoạn thẳng dốc hướng lên, lúc này đá đã tan hoàn toàn thành lỏng, toàn bộ khối lượng hỗn hợp tổng $M = 10\text{ kg}$ nước lỏng thu nhiệt nâng nhiệt độ dâng từ $0^\circ\text{C}$ đạt lên vạch mức $2^\circ\text{C}$. Khoảng thời gian nâng nhiệt chặng này: $\Delta \tau_2 = 60 - 50 = 10\text{ phút}$. Lượng nhiệt cung cấp chặng này: $Q_2 = P \cdot \Delta \tau_2 = M \cdot c \cdot (t_2 - t_1)$.
		\end{itemize}
		- Do hằng số công suất nhiệt $P$ của nguồn cung cấp đun cấp là hoàn toàn không đổi, ta lập hệ thức tỉ lệ thuận hình học giữa các lượng nhiệt lượng tiêu hao theo các khoảng chia thời gian trục hoành:
		$$\frac{Q_1}{Q_2} = \frac{P \cdot \Delta \tau_1}{P \cdot \Delta \tau_2} \Rightarrow \frac{\lambda \cdot m}{M \cdot c \cdot \Delta t} = \frac{\Delta \tau_1}{\Delta \tau_2}$$
		- Thay toàn bộ các trị số đại số số liệu đề bài vào biểu thức tỉ lệ để giải tìm khối lượng đá $m$:
		$$\frac{3,4 \cdot 10^5 \cdot m}{10 \cdot 4200 \cdot (2 - 0)} = \frac{50\text{ phút}}{10\text{ phút}} \Rightarrow \frac{340000 \cdot m}{84000} = 5$$
		$$340000 \cdot m = 5 \cdot 84000 \Rightarrow 340000 \cdot m = 420000$$
		$$\Rightarrow m = \frac{420000}{340000} \approx 1,235\text{ kg.}$$
		Vậy khối lượng nước đá ban đầu nằm trong hỗn hợp xấp xỉ bằng $1,24\text{ kg}$.
	}
\end{ex}

\Closesolutionfile{ans}

\newpage
%%%%%%%%%%%%%%%%%%%%%%%%%%%%%%%%%%%%%%%%%%%%%%%%%%%%%%%%%%%%%%%%%%%%%%%%%
% KHUNG ĐÁP ÁN BẢNG TỰ ĐỘNG IN KẾT XUẤT CỦA HỆ THỐNG
%%%%%%%%%%%%%%%%%%%%%%%%%%%%%%%%%%%%%%%%%%%%%%%%%%%%%%%%%%%%%%%%%%%%%%%%%
\begin{center} \textbf{BẢNG ĐÁP ÁN THAM KHẢO LỚP 12 --- MÃ ĐỀ \thedeso} \end{center}

\noindent \textbf{Phần I (Câu hỏi trắc nghiệm nhiều phương án lựa chọn):}